% sec-05-proposed-clinical-research.tex - Proposed Clinical Research (full-ind)
\section{Proposed clinical research}

The proposed clinical research is presented in this section of the investigational
new drug (IND) application in the form required by 21 CFR \S 312.23(a)(6). It
consists of the study protocol for the planned Phase 1 first-in-human
investigation, the informed consent that governs each participant's enrollment,
and the investigator and facilities information that establishes that the study
will be conducted by qualified persons at adequate facilities. The investigation
described here is a prospective, open-label, single-arm, first-in-human study that
combines a standard 3+3 dose-finding evaluation of perioperative daraxonrasib
(RMC-6236) with an early-feasibility evaluation of an on-premises large language
model (LLM) directed eight-arm robotic pancreaticoduodenectomy (Whipple) system, a
significant-risk Class II collaborative device. Because the device participates in
the delivery of the investigational drug, the study is conducted under a combined
IND (21 CFR \S 312) and investigational device exemption (IDE, 21 CFR \S 812)
submission with a single integrated protocol, a single informed consent carrying
the Physical AI opt-out, and a single integrated safety governance structure. The
full clinical protocol, version 1.0.0, is bound to this IND as a supporting
document and is archived with the digital object identifier
DOI 10.5281/zenodo.20780121
(\href{https://doi.org/10.5281/zenodo.20780121}{https://doi.org/10.5281/zenodo.20780121});
the synopsis and supporting tables in \S 6.1 summarize that bound document, and
every quantitative value reported here is traceable to it
\cite{daraxwhipple,onpremwhippl}.

\subsection{Study Protocol}

\subsubsection*{Protocol identification and overview}

The bound protocol is entitled the LLM-Directed PDAC Robotic Daraxonrasib Phase 1
study (protocol version 1.0.0, DOI 10.5281/zenodo.20780121). It is a prospective,
open-label, single-arm, first-in-human dose-finding and early-feasibility study in
adults with histologically or cytologically confirmed, resectable or
borderline-resectable, KRAS G12-mutated pancreatic ductal adenocarcinoma (PDAC)
who are candidates for pancreaticoduodenectomy. The central hypothesis is that the
combined drug-device intervention can be delivered with an acceptable early safety
profile and can complete the assigned operative tasks under continuous human
oversight while a recommended Phase 2 dose (RP2D) of daraxonrasib is identified,
supporting progression to a later-phase efficacy study. Enrollment is planned for
up to $n=18$ treated participants (approximately 36 screened) at a single academic
medical center, assigned across three daraxonrasib dose levels by the 3+3 rule
with staggered sentinel enrollment, at a planned accrual of approximately one to
two participants per month. The investigation is structured around the
early-feasibility pathway recognized for significant-risk devices, which generates
the device-performance and safety evidence required before any controlled
comparison is undertaken \cite{phase1guidance,daraxwhipple}.

\subsubsection*{Objectives and endpoints}

The protocol specifies three co-primary objectives, each evaluated with exact
two-sided 95\% Clopper-Pearson confidence intervals because the design is
descriptive and non-confirmatory under ICH E9 with no formal power calculation and
no multiplicity adjustment. The first co-primary endpoint is the incidence of
device- or procedure-related serious adverse events through 30 days, including any
Clavien-Dindo Grade III or higher complication. The second co-primary endpoint is
the maximum tolerated dose (MTD) and the RP2D of perioperative daraxonrasib,
determined by the 3+3 rule across a 28-day dose-limiting toxicity (DLT) window. The
third co-primary endpoint is technical feasibility, defined as the proportion of
participants completing the assigned operative tasks without unplanned conversion
attributable to device malfunction or unsafe device behavior. The secondary
endpoints are the margin-negative (R0) resection rate on day-90 pathology
(per-protocol), the International Study Group of Pancreatic Surgery (ISGPS) Grade B
or C postoperative pancreatic-fistula rate, the Clavien-Dindo Grade III or higher
complication rate at 90 days, 30-day and 90-day mortality, the conversion rate, and
operative parameters (estimated blood loss, operating-room time, and length of
stay), each benchmarked descriptively against the 2025 Dutch nationwide robotic
pancreaticoduodenectomy cohort (conversion 10.1\%, Clavien-Dindo III+ 41.3\%, ISGPS
B/C 24.4\%, 30-day mortality 3.9\%). Exploratory endpoints are progression-free and
overall survival to 24 months by RECIST version 1.1 with Kaplan-Meier estimation,
the concordance of the LLM advisory with the hard-coded safety gate, and the
sim-to-real gap in millimeters and newtons \cite{daraxwhipple,oncotrialllm}.

Table~\ref{tab:sec05-soa} presents the schedule of activities that governs the
timing of every assessment, from screening through the operative episode and the
24-month overall-survival follow-up.

\begin{table}[t]
\caption{Schedule of activities (summary) for the bound Phase 1 protocol. The
intensive safety-collection window spans consent through the 30-day postoperative
visit and encompasses the 28-day daraxonrasib dose-limiting toxicity window.}
\label{tab:sec05-soa}
\centering
\begin{tabularx}{\textwidth}{L{4.3cm} C{1.4cm} C{1.6cm} C{1.6cm} Y}
\toprule
\textbf{Assessment} & \textbf{Screen} & \textbf{Pre-op / day 0} & \textbf{Post-op to day 30} & \textbf{Long-term follow-up} \\
\midrule
Informed consent and Physical AI opt-out election & X & - & - & re-consent on amendment \\
Eligibility, ECOG, KRAS G12 genotype, staging imaging & X & - & - & - \\
Multidisciplinary resectability review & X & X & - & - \\
Daraxonrasib pre-operative course (oral once daily) & - & X & - & - \\
Protocol-mandated operative pause of daraxonrasib & - & X & - & - \\
Robotic Whipple (8 operative phases P1 to P8) & - & X & - & - \\
Pre-procedure safety matrix and USL confirmation & - & X & - & re-check after any change \\
Perioperative pause-and-restart advisory (T+7/14/21d) & - & - & X & - \\
Vitals, hematology, chemistry, hepatic/pancreatic indices & X & X & X & per follow-up schedule \\
Serum daraxonrasib trough assay & - & X & X & as indicated \\
ISGPS fistula grading (PJ dominant) & - & - & X & - \\
Adverse-event and Physical AI AE capture & X & X & X & SAE/PAI AE to 24 months \\
RECIST v1.1 imaging (central read) & X & - & X & per follow-up schedule \\
Pathology, R0 margin (day-90 read, per-protocol) & - & - & X & day 90 \\
Progression-free and overall survival & - & - & - & to 24 months (q12 weeks) \\
\bottomrule
\end{tabularx}
\end{table}

\subsubsection*{Study design and dose escalation}

The drug component proceeds under the IND pathway and the device component under
the IDE pathway, with a single integrated protocol and a single safety governance
structure. The escalation uses three dose levels anchored to the pan-KRAS RASolute
302 program, in which 300 mg once daily is the established RP2D: DL1 is 160 mg once
daily (starting dose, chosen below the reference dose to preserve a conservative
first-in-human margin), DL2 is 220 mg once daily, and DL3 is the 300 mg once-daily
reference dose. Three participants are enrolled at each level and observed across a
28-day DLT window. If no DLT occurs among the first three, the cohort escalates; if
one of three experiences a DLT, the cohort expands to six; a single DLT among six
permits escalation, whereas two or more DLTs at a level fix the MTD at the
preceding level. The MTD is the highest level at which fewer than two of six
DLT-evaluable participants experience a DLT, and the RP2D is selected at or below
the MTD. With a maximum of three levels of six, the design treats up to 18
participants. The estimation precision that this sample supports is illustrated by
the exact Clopper-Pearson intervals: an observed rate of 0 of 18 yields 0\% (95\%
CI 0 to 18\%), 1 of 18 yields 5.6\% (95\% CI 0.1 to 27\%), 0 of 3 yields 0\% (95\%
CI 0 to 71\%), and 0 of 6 yields 0\% (95\% CI 0 to 46\%) \cite{daraxwhipple}. The
analysis populations are the Safety population (all dosed participants), the
DLT-evaluable population, the Per-protocol population, and the modified
intention-to-treat (mITT) population.

The device has no pharmacologic dose; its dose analogue is the readiness evidence
required before any patient-facing activity, escalated only by passing the
mandatory Phase 0 simulation-validation gate. Three conditions must be met and
documented before the first robotic case: a Unified Safety Level (USL) rating of at
least 7.0 on the surgical-readiness scale; at least 1000 simulated procedures
completed across at least two independent simulation frameworks (for example Isaac,
MuJoCo, Gazebo, or PyBullet) so that the readiness finding is not an artifact of any
single simulator; and a sim-to-real fidelity within prespecified limits, a
trajectory gap of less than 2 mm and a tip-force gap of less than 0.5 N between the
simulation prediction and the operative record. The gate is re-run after any
software, model, or configuration change before any further patient-facing use
\cite{onpremwhippl,fdadigtwinpc}.

\subsubsection*{Population and eligibility}

The study population comprises adults with histologically or cytologically confirmed,
resectable or borderline-resectable, KRAS G12-mutated PDAC who are eligible for
perioperative daraxonrasib under the broad pan-KRAS RASolute 302 criteria and whose
tumor and vascular anatomy are compatible with the eight-arm robotic system.
Participant selection is equitable; eligibility is determined solely by the
scientific and safety criteria, and limited English proficiency is not an exclusion
criterion. The synthetic reference exemplar PAT-PDAC-0001 (a 68-year-old with a
3.4 cm head-of-pancreas tumor abutting the superior mesenteric vein at $75^\circ$,
KRAS G12D, microsatellite stable, CA 19-9 of 412 U/mL, ECOG performance status 1,
after a completed neoadjuvant modified FOLFIRINOX course of four cycles) illustrates
a representative eligible participant and does not represent an enrolled individual.
The full eligibility criteria are summarized in Table~\ref{tab:sec05-eligibility}.

\begin{table}[t]
\caption{Key inclusion and exclusion criteria for the bound Phase 1 protocol.}
\label{tab:sec05-eligibility}
\centering
\begin{tabularx}{\textwidth}{Y Y}
\toprule
\textbf{Inclusion criteria (all required)} & \textbf{Exclusion criteria (any one excludes)} \\
\midrule
Signed informed consent before any study-specific procedure, including explicit Physical AI opt-out election or declination & Distant metastatic disease on staging \\
Age 18 years or older at consent & Vascular involvement precluding a margin-negative (R0) resection, or anatomy incompatible with the eight-arm robotic approach \\
Histologically or cytologically confirmed resectable or borderline-resectable PDAC on multidisciplinary review & Prior pancreatic resection or prior surgery precluding the planned reconstruction \\
Documented KRAS G12 mutation (for example G12D, G12V, or G12R) on a validated assay & Uncontrolled intercurrent illness or comorbidity making major surgery or daraxonrasib unsafe \\
ECOG performance status 0 or 1 & Pregnancy or lactation, or unwillingness to use adequate contraception where applicable \\
Adequate organ and marrow function for major pancreaticobiliary surgery and daraxonrasib & Known contraindication or hypersensitivity to daraxonrasib or a required excipient, or an unmanageable drug interaction \\
Measured baseline CA 19-9 recorded for serial assessment; eligible for daraxonrasib under RASolute 302 pan-KRAS criteria & Inability to provide informed consent or to comply with the Schedule of Activities \\
Tumor and vascular anatomy compatible with the eight-arm robotic system on preoperative imaging & Anatomy or comorbidity otherwise incompatible with safe conduct of the combined intervention \\
\bottomrule
\end{tabularx}
\end{table}

\subsubsection*{Intervention: drug and device}

Daraxonrasib is an orally bioavailable RAS(ON) multi-selective (pan-RAS)
small-molecule inhibitor that binds the active, guanosine-triphosphate-bound state
of mutant RAS in complex with cyclophilin A, suppressing downstream effector
signaling across the common KRAS variants in PDAC; it received FDA Breakthrough
Therapy designation in June 2025 for previously treated metastatic KRAS G12-mutant
PDAC. In this study it is administered orally once daily in a perioperative
schedule comprising a defined pre-operative course, a protocol-mandated pause across
the operative window, and a postoperative restart governed by the advisory in
\S 6.1 below. The device component is an eight-arm parallel coelomic surgical
platform carrying 56 mechanical degrees of freedom (8 arms $\times$ 7 per-arm
degrees of freedom) and a 640-channel mixed-rate sensor stack (80 channels per arm),
with force channels sampling at 100 kHz and all remaining channels at the 10 kHz
command rate, and a positioning accuracy of 0.05 mm root mean square at a peak tip
velocity of 1,200 mm/s. An on-premises repository-pinned, open-weights LLM acts as a
second-opinion advisory oracle held strictly outside the robot vendor kinematic
stack and hash-pinned to a deterministic seed and a repository commit; it is never
an autonomous controller, and full autonomy is prohibited at this phase consistent
with 21 CFR \S 312.21(e). The advisory is gated by a verification-before-generation
ten-gate suite returning ACCEPT, BLOCK, or ESCALATE, aligned with H.R. 9510, the
Verification Before Generation in Physical AI Oncology Trials Act of 2026
\cite{congress9510,onpremwhippl}.

The hard cyber-physical safety envelope comprises a per-arm tip-force cap of
$\leq 3$ N, a cumulative cross-arm tip-force cap of $\leq 18$ N enforced at the
heartbeat boundary, a 10 kHz heartbeat coordination bus broadcasting a 64-byte frame
on a $100\,\mu$s per-arm watchdog deadline, an emergency arm park within $50\,\mu$s
of a missed or out-of-parameter frame, a cross-arm emergency stop (E-stop) that
halts the full platform in $\leq 3$ ms at a 0.0 N force clamp, and a
software-independent hardware E-stop meeting the 21 CFR \S 312.404 requirement of
$\leq 500$ ms. A vascular no-fly gate sampled at 10 kHz evaluates instrument-tip
proximity to five named vessels (superior mesenteric vein, portal vein, hepatic
artery, celiac axis, and superior mesenteric artery) against three concentric
thresholds: a soft-warning zone at 3.0 mm raises an LLM advisory, a no-fly zone
(1.0 mm at the superior mesenteric and portal veins; 1.5 mm at the hepatic artery,
celiac axis, and superior mesenteric artery) blocks the command, and a hard-stop
zone at 5.0 mm forces the $\leq 3$ ms E-stop; across a representative 100-tick
verdict sample the gate returned clear on 83 ticks, soft-warning on 6, no-fly on 6,
and hard-stop on 5. The operation proceeds across eight operative phases: phases P1
through P4 comprise dissection and specimen removal, phase P5 the duct-to-mucosa
pancreaticojejunostomy under a controller ring-tension band of 0.30 to 0.60 N, phase
P6 the end-to-side hepaticojejunostomy at 0.20 to 0.50 N, phase P7 the antecolic
gastrojejunostomy at 0.40 to 0.80 N, and phase P8 final imaging, hemostasis, and
drain placement \cite{onpremwhippl}.

\subsubsection*{Perioperative daraxonrasib pause-and-restart advisory}

The two components are coupled at one decision point only: the perioperative
pause-and-restart advisory for the drug is informed by the operative and pathology
results of the device procedure. Daraxonrasib is paused before surgery and restarted
postoperatively on a day chosen by an advisory that is LLM-bound and human-confirmed,
never autonomous. The advisory takes two inputs, the pancreaticojejunostomy ISGPS
fistula grade (A, B, or C) and the serum daraxonrasib trough relative to a 0.5 ng/mL
threshold, and recommends a restart at postoperative day T+7d, T+14d, or T+21d (or a
continued hold). In the 32-iteration deterministic simulation sweep, baseline serum
was 0.45 ng/mL at the operative timepoint across all iterations (below the 0.5 ng/mL
trough threshold), and the advisory recommended T+7d in 29 of 32 iterations (90.6\%,
Grade A fistula with sub-threshold trough), T+14d in 3 of 32 (9.4\%, Grade B or
delayed serum rise), and T+21d or continued hold in 0 of 32 (reserved for Grade C,
which also maps to the drug stopping rule). Figure~18 reproduces the advisory logic
and the sweep distribution.

\begin{mermaidfig}
\node[mmin]  (pause) at (0,0)        {Daraxonrasib paused\\before surgery};
\node[mmin]  (in1)   at (0,-2.0)     {Input 1: PJ ISGPS\\fistula grade A / B / C};
\node[mmin]  (in2)   at (0,-4.0)     {Input 2: serum trough\\vs 0.5 ng/mL threshold};
\node[mmproc] (llm)  at (5.6,-2.0)   {LLM advisory\\(hash-pinned)\\+ human confirmation\\(never autonomous)};
\node[mmgoal] (t7)   at (11.2,0)     {Restart T+7d\\29 of 32 (Grade A,\\sub-threshold trough)};
\node[mmstep] (t14)  at (11.2,-2.0)  {Restart T+14d\\3 of 32 (Grade B or\\delayed serum rise)};
\node[mmdark] (t21)  at (11.2,-4.0)  {Restart T+21d or hold\\0 of 32 (Grade C)};
\draw[mmedge] (pause) to[out=0,in=180,looseness=0.6] (llm.west);
\draw[mmedge] (in1) -- (llm.west);
\draw[mmedge] (in2) to[out=0,in=180,looseness=0.6] (llm.west);
\draw[mmedge] (llm) to[out=0,in=180,looseness=0.6] (t7.west);
\draw[mmedge] (llm) -- (t14.west);
\draw[mmedge] (llm) to[out=0,in=180,looseness=0.6] (t21.west);
\end{mermaidfig}
\figcaption{Figure 18. The perioperative daraxonrasib pause-and-restart advisory.
The drug is paused before surgery and the postoperative restart day is set by a
hash-pinned, human-confirmed LLM advisory (never autonomous) from two inputs: the
pancreaticojejunostomy ISGPS fistula grade (A, B, or C) and the serum daraxonrasib
trough relative to a 0.5 ng/mL threshold. In the 32-iteration deterministic sweep
the advisory recommended restart at T+7 days in 29 of 32 iterations (Grade A with
sub-threshold trough), T+14 days in 3 of 32 (Grade B or delayed serum rise), and
T+21 days or continued hold in 0 of 32 (reserved for Grade C).}

\subsubsection*{Assessments and the vascular and halt safety chain}

Clinical safety assessments comprise vital signs, a directed physical examination,
hematology, serum chemistry including hepatic and pancreatic indices, coagulation
studies, the serum daraxonrasib trough assay, and protocol-specified imaging, at the
frequencies in the schedule of activities. In addition, the device generates a
continuous safety telemetry stream that is a source of protocol observations under
21 CFR \S 312.23(g): per-arm and cumulative force profiles, trajectory accuracy
against the 0.05 mm specification, every vascular safety-zone verdict, every
heartbeat and watchdog event, and every E-stop activation, each captured at its
recorded rate in the hash-chained audit trail. The halt chain is layered: the 10 kHz
heartbeat bus broadcasts a 64-byte frame on a $100\,\mu$s per-arm watchdog deadline;
a frame that is on time continues the command state, whereas a missed or
out-of-parameter frame parks the affected arm within $50\,\mu$s and escalates to the
cross-arm E-stop that halts the platform in $\leq 3$ ms at a 0.0 N force clamp,
backed by the software-independent hardware E-stop. Postoperative pancreatic fistula
is graded by the ISGPS definition as a biochemical leak, Grade B (requiring a change
in management), or Grade C (requiring reoperation or causing organ failure or death),
and the pancreaticojejunostomy grade is the dominant determinant and the input to the
restart advisory and the drug stopping rule \cite{daraxwhipple,oncotrialllm}.

\subsubsection*{Safety reporting clocks and the six Physical AI triggers}

Adverse events are captured through two parallel and equally binding streams, a
conventional pharmacovigilance stream for clinical events graded by CTCAE version 5.0,
and a Physical AI safety stream for device- and AI-related events under 21 CFR
\S 312.32(f). Serious adverse events are reported by the investigator to the sponsor
within 24 hours of awareness, and the sponsor submits IND/IDE safety reports on the
regulatory clocks: any unexpected fatal or life-threatening suspected adverse
reaction is reported as soon as possible but no later than 7 calendar days after
initial receipt of the information, and any other serious and unexpected suspected
adverse reaction is reported no later than 15 calendar days after the determination
that the event is reportable (21 CFR \S 312.32). In addition, six Physical AI
reporting triggers are reported under \S 312.32(g): a serious Physical AI adverse
event on the 7- or 15-day timeline as applicable; a system-drug interaction causing
incorrect drug administration, regardless of seriousness, within 15 days; a
cybersecurity incident within 15 days (or 7 days if it has the potential to cause
patient harm); sustained AI model degradation below the pre-specified acceptance
criteria for three or more consecutive procedures or a cumulative 24 hours; a
sim-to-real divergence beyond twice the validated gap tolerance (trajectory deviation
greater than 4 mm or force discrepancy greater than 1.0 N); and a digital-twin
discrepancy beyond the pre-specified accuracy thresholds. For every Physical AI
adverse event, the complete hash-chained audit trail from 24 hours before the event
to 72 hours after the event is preserved under 21 CFR part 11 and made available to
the FDA on request. Figure~19 reproduces the clocks and the six triggers
\cite{phase1guidance,cfr312paiuni}.

\begin{mermaidfig}
\node[mmin]  (ev)  at (0,-2.4)      {Adverse event\\(clinical or\\Physical AI)};
\node[mmgoal] (c7) at (5.4,0)       {7 calendar days\\fatal / life-threatening\\suspected adverse reaction};
\node[mmgoal] (c15) at (5.4,-2.4)   {15 calendar days\\other serious + unexpected\\suspected adverse reaction};
\node[mmproc] (g1) at (5.4,-4.4)    {1 Serious Physical AI AE (7 / 15d)};
\node[mmproc] (g2) at (5.4,-5.8)    {2 System-drug interaction (15d)};
\node[mmproc] (g3) at (10.8,-4.4)   {3 Cybersecurity incident (15d; 7d if harm)};
\node[mmproc] (g4) at (10.8,-5.8)   {4 Model degradation ($\geq 3$ procedures / 24h)};
\node[mmproc] (g5) at (5.4,-7.2)    {5 Sim-to-real divergence $> 2\times$ ($>4$ mm / $>1.0$ N)};
\node[mmproc] (g6) at (10.8,-7.2)   {6 Digital-twin discrepancy};
\node[mmdark] (aud) at (10.8,-2.4)  {Hash-chained audit trail\\$-24$h to $+72$h\\(21 CFR part 11)};
\draw[mmedge] (ev) to[out=60,in=180,looseness=0.6] (c7.west);
\draw[mmedge] (ev) -- (c15.west);
\draw[mmedge] (ev) to[out=-80,in=180,looseness=0.5] (g1.west);
\draw[mmedge] (ev) to[out=-90,in=180,looseness=0.5] (g2.west);
\draw[mmedge] (ev) to[out=-95,in=180,looseness=0.5] (g5.west);
\draw[mmedge] (g3) -- (aud.south);
\draw[mmedge] (g4) to[out=90,in=-90,looseness=0.7] (aud.south);
\draw[mmedge] (g6) to[out=90,in=-30,looseness=0.7] (aud.south);
\end{mermaidfig}
\figcaption{Figure 19. The safety-reporting clocks and the six Physical AI reporting
triggers. A fatal or life-threatening suspected adverse reaction is reported within
7 calendar days and any other serious and unexpected suspected adverse reaction
within 15 calendar days (21 CFR \S 312.32). The six Physical AI triggers of
\S 312.32(g) are a serious Physical AI adverse event (7- or 15-day), a system-drug
interaction (15-day), a cybersecurity incident (15-day, or 7-day if patient harm is
possible), sustained model degradation (at least 3 consecutive procedures or 24
cumulative hours), a sim-to-real divergence exceeding twice the validated tolerance
(trajectory over 4 mm or force over 1.0 N), and a digital-twin discrepancy. Each
preserves the hash-chained audit trail from minus 24 to plus 72 hours under 21 CFR
part 11.}

\subsubsection*{Discontinuation and stopping rules}

The protocol distinguishes discontinuation of an intervention from withdrawal from
the study; a participant who stops one or both components for a protocol-defined
reason remains in the study for safety and survival follow-up unless consent for that
follow-up is withdrawn. The investigational device is discontinued for the index
procedure, with conversion to the appropriate fallback technique, whenever any of
five prespecified conditions occurs: a conversion to open or manual technique judged
by the supervising surgeon to better serve patient safety or oncologic adequacy
(conversion is always available and is never itself a protocol violation); repeated
excursions against the $\leq 3$ N per-arm or $\leq 18$ N cumulative caps; any breach
of a vascular no-fly zone (1.0 mm at the superior mesenteric and portal veins;
1.5 mm at the hepatic artery, celiac axis, and superior mesenteric artery) or a
recurring pattern of no-fly approaches; any sim-to-real divergence beyond twice the
validated tolerance (trajectory deviation greater than 4 mm or force discrepancy
greater than 1.0 N), which also triggers a \S 312.32(g) report; and any hard-stop
cross-arm E-stop event, after which device resumption requires that the cause be
identified and the system re-pass the pre-procedure safety matrix. Daraxonrasib is
discontinued in an individual participant for a drug-attributable DLT during the
28-day window, for any other unacceptable toxicity not resolving with protocol dose
interruption and reduction, for a confirmed ISGPS Grade C fistula (which precludes
restart and maps to the T+21d or continued-hold arm of the advisory), for pregnancy,
for contraindicating intercurrent illness, or at the participant's request
\cite{onpremwhippl,daraxwhipple}.

Two binding halt rules govern the study as a whole and are applied by the data
safety monitoring board (DSMB): a device- or procedure-related serious adverse event
in at least one of three participants within a cohort, or two device-related deaths
at any time, triggers DSMB review and a halt. On any hold or closure, enrollment and
patient-facing robotic activity stop immediately, affected participants are
transitioned to appropriate clinical care, and a controlled decommissioning and
data-preservation procedure seals the complete hash-chained command and telemetry
record for the regulatory retention period.

\subsubsection*{Statistical considerations and oversight}

The statistical approach is descriptive and estimation-based, non-confirmatory, and
follows ICH E9 and the TRIPOD+AI reporting standard for the model-assisted advisory
component. No formal power calculation is performed; exact two-sided 95\%
Clopper-Pearson confidence intervals are reported for the binomial endpoints, and a
two-sided alpha of 0.05 is used only in a supportive sense with no multiplicity
adjustment. Independent safety oversight operates on three tiers: the DSMB reviews
accumulating safety data at each cohort boundary and at any triggered review and
holds halt authority; the Independent Safety Monitor (ISM) is designated for each
procedure and provides real-time case-level oversight with the authority to call a
stop; and the Physical AI Safety Review Committee meets on a cadence not exceeding
90 days (and ad hoc on any triggering event) to review the autonomy-bearing behavior
of the system, the verification-before-generation gate-surface decisions, the USL
rating, and any software change, and to recommend USL reassessment or oversight
changes. Seven Physical AI record types (the LLM advisory record, the robot command
record, the sensor and telemetry record, the safety-limit-event record, the
verification-before-generation gate-surface record, the human-override and
intervention record, and the binding hash-chained audit trail) are retained for the
period required by 21 CFR \S 312.57, at least two years after a marketing
application is approved or, if none is filed, two years after the investigation is
discontinued and the FDA is notified \cite{daraxwhipple,clintrustpai}.

The full protocol, version 1.0.0, with its complete schedule of activities,
statistical analysis plan, and Physical AI System Description meeting the eleven
fields of 21 CFR \S 312.23(g), is the bound supporting document for this section and
is archived at DOI 10.5281/zenodo.20780121
(\href{https://doi.org/10.5281/zenodo.20780121}{https://doi.org/10.5281/zenodo.20780121}).

\subsection{Informed Consent}

Informed consent is obtained by a qualified member of the research team who is not
the treating surgeon, before any study-specific procedure, in language the
participant can understand and with adequate time and opportunity to ask questions,
consistent with 21 CFR part 50 and ICH E6(R3) \cite{cfr050paiuni,ichpaistduni}. The
consent process incorporates each of the basic and additional elements of informed
consent required by 21 CFR \S 50.25, and a copy of the form and any participant
information sheet are submitted with the protocol. The consent discussion covers the
diagnosis, the alternatives (including standard open or conventional robotic
pancreaticoduodenectomy without the investigational advisory layer), and the
prognosis of a malignancy with a total five-year survival below 13\%; the
perioperative daraxonrasib regimen, its once-daily oral administration, the
protocol-mandated operative pause, and the toxicities characteristic of the RAS(ON)
inhibitor class; and, in addition to the standard elements, an explicit and distinct
disclosure of the on-premises LLM-directed eight-arm robotic system, its role as a
second-opinion advisory oracle held outside the vendor kinematic stack, its
continuous human oversight with immediate manual override, and its system-specific
risks. The voluntary nature of participation, the right to withdraw at any time
without penalty and without loss of any benefit to which the participant is otherwise
entitled, and the maintenance of confidentiality under the HIPAA Safe Harbor
de-identification method are stated at consent and reaffirmed at each contact.

The consent form addresses the therapeutic misconception directly. Because the
population is vulnerable, facing the most technically demanding curative-intent
abdominal operation for a lethal cancer, the form and the consent conversation make
explicit that the study is a first-in-human, dose-finding, and early-feasibility
investigation whose primary purpose is to characterize safety and feasibility rather
than to confer individual therapeutic benefit, and that the LLM directs no autonomous
action. The staged-autonomy model is explained: in this Phase 1 the system operates
only at the clinician-controlled (teleoperation or supervised) and supervised
semiautonomous levels under continuous human oversight, full autonomy is prohibited
consistent with 21 CFR \S 312.21(e), and any future move to higher autonomy is
reserved for a later phase and gated by independent review. An independent consent
monitor is used when indicated, and qualified interpreters and translated materials
are provided so that consent is meaningful for participants with limited English
proficiency \cite{clintrustpai,oncotrialllm}.

A defining feature of this trial class is the documented Physical AI opt-out. Under
21 CFR \S 312.60(f), as adapted for Physical AI investigations, the participant is
offered, and may exercise, a documented right to request administration of the
operative intervention without the Physical AI advisory layer where clinically
feasible, without penalty and without affecting eligibility for the remainder of the
study or for clinical care \cite{cfr312paiuni}. The participant's election or
declination of the opt-out is recorded in the source document and the case report
form, and the schedule of activities (Table~\ref{tab:sec05-soa}) captures that
election at the consent visit. Where an unanticipated problem changes the risk
profile of the combined intervention, including a relevant Physical AI safety event,
the consent document is updated and currently enrolled participants are re-consented
before any further study procedure, consistent with the original opt-out.

\subsection{Investigator and Facilities Data}

The investigation is conducted at a single academic medical center through its
institutional multidisciplinary pancreatic-cancer program. Eligible individuals are
identified at tumor-board review and at surgical and medical-oncology clinics, and
the study is described to them by a member of the research team who is not the
treating surgeon, to reduce the potential for undue influence. The site possesses an
established high-volume pancreaticobiliary surgical service, a robotic surgical
program, an investigational-product pharmacy with controlled, temperature-monitored
storage and a drug-accountability log, an anatomic-pathology service applying a
standardized synoptic margin protocol with masked independent adjudication, a central
imaging capability for RECIST version 1.1 reads by a blinded reviewer, and the
clinical laboratory capacity for the serum daraxonrasib trough assay and the
hematologic, hepatic, renal, and pancreatic safety panels. The device facilities
include the qualified clinical environment in which the eight-arm platform is
maintained and calibrated between procedures, the operator control stations, the
on-premises compute environment that hosts the repository-pinned advisory LLM in
network isolation from the vendor kinematic stack under a deny-by-default
authentication and authorization model, and the 21 CFR part 11 records environment
that maintains the hash-chained, append-only audit trail. Day-to-day monitoring is
delegated to a robotics- and artificial-intelligence-competent clinical research
organization (CRO) over the site team, which comprises the operating surgeon, a
primary operator, and a qualified backup operator ready to assume control or invoke
the manual override at any time \cite{onpremwhippl,paisitedocpk}.

The governance structure that supervises these investigators and facilities is the
three-tier oversight architecture described in \S 6.1: the Sponsor-Investigator
answers to a reviewing IRB and to the independent DSMB, convenes the Physical AI
Safety Review Committee on a cadence not exceeding 90 days, designates the
Independent Safety Monitor for each procedure, and delegates monitoring to the CRO.
Figure~20 presents the governance and safety-oversight structure that binds the
investigators, the facilities, and the independent reviewers together.

\begin{mermaidfig}
\node[mmgoal] (si)  at (4.2,0)       {Sponsor-Investigator\\ChemicalQDevice};
\node[mmstep] (irb) at (0,-2.6)      {Reviewing IRB\\SR device + IND};
\node[mmstep] (dsmb) at (4.2,-2.6)   {Data and Safety\\Monitoring Board\\(halt authority)};
\node[mmstep] (pais) at (8.4,-2.6)   {Physical AI Safety\\Review Committee\\$\leq 90$-day cadence};
\node[mmin]  (cro)  at (-1.0,-5.2)   {CRO / clinical\\monitor (robotics\\+ AI competent)};
\node[mmstep] (ism) at (8.4,-5.2)    {Independent Safety\\Monitor (per\\procedure, stop authority)};
\node[mmin]  (site) at (4.2,-7.4)    {Site team: surgeon,\\operator, backup\\operator};
\draw[mmedge] (si) -- (irb.north);
\draw[mmedge] (si) -- (dsmb);
\draw[mmedge] (si) -- (pais.north);
\draw[mmedge] (si) to[out=180,in=90,looseness=0.7] (cro.north);
\draw[mmedge] (cro) to[out=-90,in=180,looseness=0.7] (site.west);
\draw[mmedged] (dsmb) -- node[mmlabel,pos=0.55]{stop / continue} (site.north);
\draw[mmedged] (pais) -- node[mmlabel,pos=0.5]{USL changes} (ism.north);
\draw[mmedged] (ism) to[out=-90,in=0,looseness=0.7] (site.east);
\end{mermaidfig}
\figcaption{Figure 20. Study governance and safety oversight. The
Sponsor-Investigator answers to the reviewing IRB and the DSMB (halt authority),
convenes a Physical AI Safety Review Committee on a cadence not exceeding 90 days,
designates an Independent Safety Monitor (per-procedure, stop authority), and
delegates monitoring to a robotics- and AI-competent CRO over the site team (surgeon,
operator, and a qualified backup operator).}

\subsubsection{FDA Form 1572}

A Statement of Investigator (FDA Form 1572) is executed by each investigator
responsible for the conduct of the study at the clinical site, in satisfaction of
21 CFR \S 312.53(c), and is assembled with the submission together with the
investigator's current curriculum vitae and the documentation of relevant
qualifications. The executed Form 1572 attests that the investigator has the
training and experience to assume responsibility for the proper conduct of the
investigation, identifies the name and address of the clinical site and of the
clinical laboratory and other facilities to be used, names the reviewing IRB
responsible for review and approval of the study, lists the sub-investigators (the
operating surgeon, the primary and backup device operators, the medical oncologist,
the pathologist, and the research coordinators) who will assist in the conduct of
the investigation, and identifies the protocol by its title and version. By signing,
the investigator commits to conduct the study in accordance with the protocol, to
personally conduct or supervise the investigation, to inform participants that the
drug and device are being used for investigational purposes and to ensure that
consent and IRB requirements are met, to report adverse events to the sponsor in
accordance with 21 CFR \S 312.64, to read and understand the investigator's brochure
and the Physical AI System Description, to ensure that all associates and staff are
informed of their obligations, to maintain adequate and accurate records available
for inspection under 21 CFR \S 312.68, and to comply with all other requirements of
21 CFR part 312. The executed forms, with attachments, are bound into the submission
and are updated whenever a sub-investigator or facility changes \cite{phase1guidance}.

\subsubsection{Disclosure of Financial Interests}

Financial interests are disclosed and managed in accordance with 21 CFR part 54,
Financial Disclosure by Clinical Investigators, and the certification and disclosure
statements (FDA Forms 3454 and 3455 as applicable) are collected before participation
and updated through one year after completion of the study. The disclosure makes
explicit the principal conflict of interest in this investigation: the
Sponsor-Investigator, Kevin Kawchak, is the chief executive officer of
ChemicalQDevice, the entity that develops the on-premises LLM-directed eight-arm
robotic Physical AI system under study and holds the combined IND and IDE. This
relationship is disclosed in full to the FDA and to the reviewing IRB, and it is
managed by the independent oversight structure described in \S 6.3 so that safety and
escalation decisions are not made by a party with a financial interest in the outcome
alone \cite{cfr312paiuni,clintrustpai}.

The mitigation is structural and operates at every decision point that could be
influenced by the conflict. The DSMB, composed of members independent of the conduct
of the study with surgical, oncologic, biostatistical, and device-safety expertise,
holds the binding halt authority and applies the prespecified halt rules (a device-
or procedure-related serious adverse event in at least one of three participants
within a cohort, or two device-related deaths at any time). The Independent Safety
Monitor, designated for each procedure and independent of the operating team, holds
real-time stop authority at the case level. The Physical AI Safety Review Committee
reviews the autonomy-bearing behavior of the system, the verification-before-generation
gate-surface decisions, the USL rating, and any software change on a cadence not
exceeding 90 days, and recommends USL reassessment or oversight changes. Routine
monitoring is delegated to a CRO independent of the device manufacturer, and audits
may be conducted by a designee independent of the conduct of the study. Because the
binding safety, escalation, and halt decisions rest with these independent bodies
rather than with the Sponsor-Investigator, the financial conflict is disclosed and
mitigated to a degree commensurate with the first-in-human, autonomy-bearing nature
of the investigation, and any identified conflict is disclosed and managed before it
can affect the conduct, analysis, or reporting of the study \cite{daraxwhipple}.
